\documentclass[11pt]{article}
\usepackage[a4paper,margin=25mm]{geometry}
\usepackage{amsmath,amssymb,amsthm,hyperref}
\emergencystretch=2em
\newtheorem{theorem}{Theorem}\newtheorem{lemma}[theorem]{Lemma}\newtheorem{proposition}[theorem]{Proposition}
\title{Boundary constraints for equal-area rectangle dissections}
\author{Private recovery manuscript for L018}\date{12 October 2026}
\begin{document}\maketitle
A perfect partition in this paper is an axis-parallel dissection of a rectangle
into at least two pairwise noncongruent rectangles of equal positive area.
Congruence allows a quarter-turn. This paper supplies complete versions of the
boundary and search arguments used by Garc\'ia-Col\'in, Leemans, M\"u\ss ig and
Rold\'an~\cite{source}; the finite exclusions themselves require separately
completed computations. No universal nonexistence assertion follows here.

\begin{theorem}[Source Lemma 2.1]\label{seven}
A perfect partition has at least seven pieces, even if their side lengths are real.
\end{theorem}
\begin{proof}
Suppose otherwise and choose such a partition with the least number $r$ of pieces.
There cannot be two pieces: they would have equal widths or equal heights and,
by equality of area, would be congruent. No piece spans an entire width or height
of the container. Removing such a piece leaves one or two nonempty rectangular
regions. A region consisting of one piece would make that piece congruent to
the removed spanning piece. Every remaining region therefore contains at least
two pieces and is a smaller perfect partition, contradicting minimality. Each container side
therefore meets at least two pieces along intervals of positive length, and the
four corner pieces are distinct.

Suppose each side meets exactly two pieces. These must be its corner pieces, so
only the four corner pieces meet the boundary. Scale the two coordinate axes
independently to make the container the unit square. This scaling is used only
to derive equality of the widths and heights of two pieces, which remains an
equality on undoing the scaling. The common normalized area is $a=1/r$. Let $x$
and $y$ be the widths of the bottom-left and top-left pieces. The widths on the
right are $1-x$ and $1-y$. Their heights and the two vertical boundary equations
give
\[
 y=\frac{ax}{x-a},\qquad 1-y=\frac{a(1-x)}{1-x-a}.
\]
The denominators are positive because the two lower pieces have heights strictly
less than one. Adding the equations and clearing denominators gives
$(1-2a)x(1-x)=a(1-2a)$. Since $r>2$, this implies $x(1-x)=a$ and $y=1-x$.
The bottom-left and top-right pieces thus have the same width $x$ and the same
height $a/x$, contrary to noncongruence.

Some side consequently meets $k\ge3$ pieces. Take this as the bottom side.
The heights of these pieces are pairwise distinct: equal heights and equal
areas would force congruence. None reaches the top. Each has at least one piece
immediately above its upper edge, and these sets of upper neighbors are mutually
disjoint, since an upper neighbor's bottom coordinate is the height of the
corresponding bottom piece. They also contain no bottom piece. The lowest
bottom piece has at least two upper neighbors. Otherwise its unique upper
neighbor covers its full top edge. At each endpoint, either the container wall
or the adjacent, strictly taller bottom piece prevents horizontal extension,
so this neighbor has exactly the same width. Equal area again gives congruence.
There are therefore at least $k+(k+1)\ge7$ pieces.
\end{proof}

\begin{lemma}[Integral positions]\label{integral}
In an axis-parallel tiling with integer piece widths and heights, every piece
corner has integer coordinates relative to a container corner.
\end{lemma}
\begin{proof}
Choose a horizontal line through the interior of a particular vertical piece
edge, avoiding all horizontal piece edges. The segments into which the tiling
partitions this line have the integer widths of the pieces it crosses.
Every endpoint's horizontal coordinate is their partial sum from the left wall.
A vertical line gives the analogous conclusion for horizontal edges.
\end{proof}

\begin{proposition}[Exact finite search interface]\label{search}
Fix positive integers $W,H,r$ with $r\mid WH$, and set $A=WH/r$.
Let the available shapes be the unordered pairs $\{a,b\}$ with $ab=A$ and
at least one orientation fitting the container. A perfect partition into $r$
pieces exists if and only if the following finite skyline search succeeds.
Initially all $W$ unit columns have height zero. At each step choose the leftmost
column of minimum height $y$, with horizontal coordinate $x$. Try every unused
shape in each fitting orientation $w\times h$ for which the full interval
$[x,x+w)$ has current height $y$ and $y+h\le H$. Raise that interval by $h$ and
continue. Success means $r$ pieces have been placed; otherwise exhaust the choices.
\end{proposition}
\begin{proof}
Every placement is inside the box and disjoint from earlier placements, and no
unordered shape is reused. After $r$ placements the total area is $rA=WH$, so
there is no uncovered open region. Conversely, take an actual partition and
follow the prescribed order. By Lemma~\ref{integral}, all its coordinates are
integer. The piece covering the leftmost lowest uncovered unit cell must begin
at $(x,y)$: its bottom cannot lie below the minimum uncovered level, and extending
to the left at that level would intersect the covered region or leave the box.
The columns across its width must all have the same height $y$, since a higher
one would intersect its interior and a lower one contradicts minimality.
Its exact orientation is one of the search choices. This inductively recovers
the actual partition. The search has finite depth $r$ and finite branching.
\end{proof}

\begin{lemma}[Choice of uncovered valley]\label{valley}
The exact search may instead choose any maximal flat skyline interval whose
neighbors have strictly larger covered heights, with container walls allowed
as neighbors. In particular it may choose such an interval with fewest
currently fitting shapes and orientations.
\end{lemma}
\begin{proof}
Let the interval be $[x,x+w)$ at height $y$. In a completion, the piece
covering its leftmost lowest uncovered unit cell cannot begin below $y$,
since that would overlap the covered region. It cannot extend to the left
because of the wall or the strictly higher neighboring column. It therefore
begins at $(x,y)$. It cannot cross the strictly higher right neighbor either.
Thus its entire base lies in this flat interval and is one of the legal
choices. Each placement still preserves the skyline invariant. Induction on
the remaining pieces proves completeness for any deterministic selection of
such an interval. A lowest interval always exists until completion.
\end{proof}

\begin{proposition}[Safe necessary prunings]\label{pruning}
The following conditions may be required at every partial skyline state without
losing any completion:
\begin{enumerate}
\item each remaining open component separated by full-height columns has area
a multiple of $A$;
\item the width of every lowest uncovered plateau is the sum of widths of a
subset of distinct unused shapes in fitting orientations;
\item the four container sides admit side subsets with their exact lengths;
adjacent subsets share exactly their common corner piece, and opposite subsets
can share only pieces spanning the intervening container dimension.
\end{enumerate}
For a square, one may additionally require that its bottom-left corner piece
has least unordered shape index among all four corner pieces, and that this
piece has width at most height. Repeated corner identities must be allowed.
\end{proposition}
\begin{proof}
An unplaced rectangle cannot cross a full-height column, so the first assertion
follows by summing equal areas within a component. The pieces starting along a
lowest plateau partition its base interval and have distinct unused shapes,
proving the second assertion. The third follows by reading the boundary of an
actual tiling; a rectangle meeting two opposite sides spans the intervening
width or height. Finally the dihedral symmetries of a square can send a least
corner piece to the bottom left. Reflection in the diagonal through that corner
then chooses its width at most its height. A piece spanning a dimension may
occupy two corners, which is why the least-index inequalities are non-strict.
\end{proof}

\begin{proposition}[Parity endpoint constraint]\label{parity}
In any integer tiling, take the parity of the number of uncovered unit cells on
each horizontal row of a partial placement. Every remaining rectangle of odd
width contributes the indicator of its vertical interval, and rectangles of even
width contribute zero. Thus the discrete derivative of the uncovered parity
profile is the symmetric difference of the endpoint pairs of the remaining
odd-width intervals. The analogous assertion holds for columns and odd heights.
Also, the alternating column sum of the uncovered cells is the sum of
$\pm h$ over remaining odd-width pieces, and the alternating row sum is the
sum of $\pm w$ over remaining odd-height pieces.
\end{proposition}
\begin{proof}
Reduce the width contributions on a row modulo two and then take differences
between consecutive rows, including the two exterior zero rows. An interval's
indicator changes exactly at its two endpoints. For the alternating sum, the
sum of $(-1)^x$ along an integer interval is zero for even length and is $1$ or
$-1$ for odd length. Sum over the height of each rectangle. Columns are dual.
\end{proof}

\begin{proposition}[A bounded exact reduction for the hard square]
In a perfect partition of the $840$-square into $21$ pieces, a quarter-turn
permits the four odd-sided pieces to have widths $75,105,175,525$ and respective
heights $64(7,5,3,1)$. Their bottom coordinates have the form $t+64k_i$, where
$0\le t$, and the normalized tuple $(k_1,k_2,k_3,k_4)$ is one of
\[
\begin{gathered}
 (0,0,4,4),\quad (0,0,5,7),\quad (0,2,0,2),\\
 (0,3,0,7),\quad (1,0,5,0),\quad (1,3,0,0).
\end{gathered}
\]
This necessary reduction is not an assertion that any such partition exists.
\end{proposition}
\begin{proof}
The smaller factors of $33600$ whose partners fit the $840$-square are
\[
\begin{gathered}
40,42,48,50,56,60,64,70,75,80,84,\\
96,100,105,112,120,140,150,160,168,175.
\end{gathered}
\]
Thus there are exactly $21$ shapes of area $840^2/21$, and all must be used.
The four with an odd side are
\[
 (75,448),\quad (105,320),\quad (175,192),\quad (525,64).
\]
All others have two even sides.
By the alternating-sum part of Proposition~\ref{parity}, the even partners of
those oriented with odd width must have a signed sum zero, and the same holds
for those oriented with odd height. After division by $64$, the partners are
$7,5,3,1$. No nonempty proper subset has a signed zero sum: pairs have unequal
members, odd-cardinality subsets have odd sum, and singletons are nonzero.
Consequently the four odd dimensions lie along a common axis. Rotate if needed.

The entire square has an even number of cells on every row. Form a multigraph
whose vertices are the endpoints of the four odd-width pieces' vertical
intervals, and whose edges are those intervals. Every vertex has even degree
by Proposition~\ref{parity}. There are no loops. A two-edge component would
have equal edge lengths, which the four distinct heights exclude. Hence this
four-edge graph is connected and is a four-cycle; any repeated vertex in a
closed traversal would split it into two forbidden two-edge cycles. All edge
lengths are multiples of $64$, so its vertices share one residue modulo $64$.
The unique partition of $\{7,5,3,1\}$ into equal sums is
$7+1=5+3$. Start a traversal with the positively directed length-$7$ edge;
the remaining signed lengths are $1,-5,-3$ in one of their six orders.
The six resulting interval-start tuples, after subtracting their minimum,
are precisely those displayed. The alternating-column equation additionally
forces the starting $x$ parities of pieces 1 and 4 to agree, and those of
pieces 2 and 3 to be the opposite parity: these are the two sides of
$7+1=5+3$. This is an elementary six-case calculation,
not a packing census. All piece heights in the chosen orientation are even,
so the integral-position argument also implies that $t$ is even.
\end{proof}

\begin{proposition}[The other hard square]\label{twenty}
In an $840$-square perfect partition with $20$ pieces, exactly eight of the
nine odd-sided shapes are selected. Their orientation and relevant starting
coordinate parity belong to the following finite, explicitly specified list.
Put
\[
 (w_0,\ldots,w_8)=(49,3,45,35,5,21,7,9,15).
\]
Choose one omitted index $o$, a set $X\subseteq\{0,\ldots,8\}\setminus\{o\}$,
and subsets $P_X\subseteq X$, $P_Y\subseteq Y$, where
$Y=\{0,\ldots,8\}\setminus(\{o\}\cup X)$, subject to
\[
 2\sum_{i\in P_X}w_i=\sum_{i\in X}w_i,\qquad
 2\sum_{i\in P_Y}w_i=\sum_{i\in Y}w_i.                 \tag{*}
\]
The pieces in $X$ have odd width and start at $x$ of parity $1$ exactly for
indices in $P_X$; the pieces in $Y$ have odd height and start at $y$ of parity
$1$ exactly for indices in $P_Y$. There are exactly $44$ such configurations.
A quarter-turn allows the additional condition $|X|\ge4$, leaving $26$.
These counts are exhaustible by the displayed nine-bit subset definition;
the exact program reconstructs that definition rather than assuming a stored
configuration list.
\end{proposition}
\begin{proof}
There are exactly $22$ fitting unordered shapes of area $35280$. The nine
with an odd side, in the indicated order, are
\[
\begin{gathered}
(45,784),(48,735),(49,720),(63,560),(80,441),\\
(105,336),(112,315),(144,245),(147,240).
\end{gathered}
\]
Their even partners divided by $16$ are precisely the odd numbers $w_i$.
All other shapes have both sides even. Since $20$ of $22$ shapes are used,
at least seven of these nine odd-sided shapes must be used. For either axis,
the corresponding alternating-sum equation of Proposition~\ref{parity}, divided
by $16$ and reduced modulo two, says that the number of selected odd sides
on that axis is even. The total selected number is therefore even and hence
is exactly eight. The two exact alternating sums are precisely (*): changing
all signs has no effect on zero. These necessary equations exhaust the asserted
finite list by considering every omitted index and every subset, with no
packing inference from a mere compatible configuration.

A quarter-turn exchanges the roles of $X$ and $Y$ (possibly complementing
one parity subset, which preserves (*)), so at least four can be horizontal.
Reflections in either container midline preserve this orientation condition
and preserve the list: a reflection along the odd dimension complements its
parity subset and leaves a zero signed sum zero. These two reflections act
transitively on the four corners. Thus one may still put a least-index corner
shape at the bottom left, using non-strict comparisons if corners coincide.
One may \emph{not} also require that corner's width to be at most its height.
Exactly this last restriction is omitted in both specialized $840$ searches.
The same reflection argument applies to the all-odd-width normalization for
$21$ pieces. This establishes all symmetry reductions used in those programs.
\end{proof}

\paragraph{Whole-program and source boundaries.}
The independent C++ implementations use a
run-length skyline, explicit unordered identities, integer arithmetic, and the
stated necessary prunings; the valley and two specialized versions add only
the proved valley order and parity restrictions. A deadline produces \texttt{UNKNOWN}, never an
exclusion. Its source must be read and a completed exhaustive run must be bound
before any finite exclusion is accepted. The independent boundary relaxation
can prove infeasibility only; a feasible boundary assignment is not a tiling.
The linked repository edition retrieved at commit
\texttt{5af96dad3646c45e2e0cf9d62481b256c22569a0} is not asserted to be the
historical execution revision. In its rectangular routines, a filtered list of
orientations need not be closed under reversal: thus dividing its length by two
and identifying its reverse index with a rotated piece are not valid in general.
This is a current-edition code qualification, not a counterexample to the source's
bounded nonexistence assertions. Their complete new finite recovery is supplied in the companion application;
scientific adoption remains separate.

\begin{thebibliography}{1}
\bibitem{source} N. Garc\'ia-Col\'in, D. Leemans, M. M\"u\ss ig and \'E. Rold\'an,
\emph{There is no perfect Mondrian partition for squares of side lengths less
than 1001}, arXiv:2311.02385v1, eleven-page fixed source edition (2023).
\end{thebibliography}
\end{document}
